\documentclass[11pt,a4paper]{article}
\usepackage{amsmath,amssymb,amsthm,bm}
\usepackage{geometry}
\geometry{margin=1in}
\usepackage{hyperref}
\usepackage{graphicx}
\usepackage{setspace}
\usepackage{booktabs}
\onehalfspacing

\title{\textbf{Paper OOL-11: Photochemical Energy Capture and Overload Stabilization:\\ Functional Bioenergetic Materials After the Core Gate}}
\author{Steve Bico Mujjabi, MD\\
\small Independent Researcher\\
\small ORCID: \href{https://orcid.org/0009-0001-0556-5516}{0009-0001-0556-5516}}
\date{May 2026}

\begin{document}
\maketitle

\begin{abstract}
Papers 1--10 are sufficient to describe a geochemical and protocellular path toward primitive life without requiring modern photosynthesis. This manuscript adds a later expansion channel: photochemical energy capture. The corrected claim is functional rather than molecule-specific. ``Chlorophyll-like'' means energy-input amplification, not literal chlorophyll at the earliest origin; ``melanin-like'' means surplus stabilization, not a required first pigment. CDFD predicts that once usable input rises, systems must couple capture to electron transport, proton/ion transport, retention, closure, information stability, and overload buffering. The paper therefore preserves the plant and pigment insights as mature exemplars while keeping the origin-of-life mechanism chemically open.
\end{abstract}

\section{Why This Paper Exists}
Papers 1--10 build the core chain: thermodynamic framing, geochemical redox/proton architecture, oscillating polymerization, autocatalytic closure, adaptive boundaries, replication/protocell integration, and parasitic boundary pressure. That sequence can describe chemolithotrophic or geochemical routes toward proto-life~\cite{Russell1994,Martin2007,Lane2010}.

The missing question is different: what happens when a system gains access to a stronger or more spatially independent energy source, especially light? Photochemical input can raise the available driving flux without requiring permanent attachment to one vent geometry. However, naming modern chlorophyll too early makes the theory historically brittle. The corrected principle is:
\begin{quote}
CDFD predicts functional roles, not mandatory molecules.
\end{quote}

\section{Functional Material Classes}
\begin{center}
\small
\begin{tabular}{p{0.21\textwidth}p{0.15\textwidth}p{0.22\textwidth}p{0.31\textwidth}}
\toprule
CDFD role & Symbol & Modern example & Prebiotic candidate class \\
\midrule
Energy-input amplification & $I_E$ & Chlorophyll antenna & Fe--S photoredox, porphyrinoids, flavins, quinones, mineral semiconductors \\
Electron transport & $\sigma_e$ & ETC cofactors, cytochromes & Fe--S clusters, magnetite-like phases, mixed-valence minerals \\
Proton/ion coupling & $\sigma_p$ & Proton motive force & Interfacial water, pores, pH gradients, mineral-water boundaries \\
Retention & $\Xi_{\mathrm{loc}}$ & Membranes, organelles & Mineral pores, gels, coacervates, LLPS droplets \\
Overload stabilization & $B_{\mathrm{stab}}$ & Carotenoids, NPQ, melanins & Polyphenols, disordered aromatics, radical sinks, mineral screens \\
\bottomrule
\end{tabular}
\end{center}
The names are mnemonic labels for functional channels. They should not be read as claims that those exact substances were present at the earliest transition.

\section{Capture Before Stabilization}
Let $I_E$ denote usable environmental energy input. A minimal throughput expression is
\begin{equation}
    J_{\mathrm{bio}}=I_E\frac{\sigma_e\sigma_p}{B_{\mathrm{stab}}}.
    \label{eq:functional-throughput}
\end{equation}
Here $B_{\mathrm{stab}}$ is a net stabilization burden: radical load, thermal disorder, photodamage risk, hydrolysis pressure, and uncoupled side reactions. A stabilizing material lowers effective $B_{\mathrm{stab}}$ only if it absorbs or redirects excess flux without blocking productive electron/proton coupling.

The corrected causal order is therefore
\begin{equation}
\text{capture}\rightarrow\text{couple}\rightarrow\text{retain}
\rightarrow\text{close}\rightarrow\text{stabilize surplus}.
\end{equation}
Chlorophyll-like functions belong at the capture stage. Melanin-like functions belong at the surplus-stabilization stage.

\section{Pre-Chlorophyll Photochemical Amplification}
Modern chlorophyll is a mature endpoint, not the primitive assumption. Earlier energy-input materials could have performed weaker versions of the same role:
\begin{itemize}
    \item Fe--S and metal sulfide surfaces that couple light or redox energy to surface reactions,
    \item semiconducting minerals that create electron-hole separation under illumination,
    \item porphyrin-like or metal-organic precursors with conjugated absorption,
    \item flavins, quinones, and aromatic cofactors that shuttle electrons,
    \item retinal-like or other simple chromophore systems in later protocellular settings.
\end{itemize}
The historical question is empirical: did any candidate raise $I_E$ while remaining coupled to $\sigma_e$, $\sigma_p$, and retention?

\section{Overload Stabilization}
Once photon capture or chemical input becomes strong, excess flux becomes dangerous. Productive throughput is limited by the slowest coupled channel:
\begin{equation}
    J_{\max}\leq \min(I_E,\sigma_e,\sigma_p,1/B_{\mathrm{stab}}).
\end{equation}
If $I_E$ exceeds the capacity of electron transport, proton/ion coupling, carbon fixation, repair, or dissipation, the surplus becomes damage:
\begin{equation}
    \Psi_{s,\mathrm{overload}}
    =
    \frac{I_E}{C_{\mathrm{use}}+C_{\mathrm{diss}}}.
\end{equation}
Modern plants solve this with carotenoids, xanthophyll cycling, non-photochemical quenching, repair of PSII, and related dissipative controls. Melanin-like polymers and other disordered aromatics are mature examples of broadband buffering and redox smoothing, not required first molecules~\cite{Cleaves2008,Higgs2015}.

\section{Photochemical Fe--S Bridge}
Photochemistry need not be late or plant-only. UV-driven prebiotic synthesis of Fe--S clusters suggests that light can participate in forming redox cofactors or cofactor-like clusters before modern photosystems~\cite{Bonfio2017}. The corrected order is therefore not
\[
\text{dark geochemistry}\rightarrow\text{chlorophyll}\rightarrow\text{plants}.
\]
It is
\[
\text{redox geochemistry}
\leftrightarrow
\text{photoredox chemistry}
\rightarrow
\text{mature photosynthetic pigments}.
\]

\section{Numerical Visualization}
\begin{figure}[ht]
\centering
\includegraphics[width=\textwidth]{outputs/paper_11/photochemical_capture_before_buffering.png}
\caption{Photochemical capture before surplus stabilization. The first panel treats chlorophyll as the mature endpoint of a wider energy-capture function. The second panel shows why buffering becomes relevant only after captured input creates overload.}
\end{figure}

\section{Falsifiability}
The corrected model separates functional claims from material claims:
\begin{center}
\small
\begin{tabular}{p{0.18\textwidth}p{0.36\textwidth}p{0.36\textwidth}}
\toprule
Layer & Claim & Falsification route \\
\midrule
Functional & Proto-life requires sustained input above maintenance & Demonstrate self-maintaining heredity without persistent throughput \\
Coupling & Electron and proton/ion channels must be co-localized & Demonstrate robust proto-metabolism with uncoupled charge channels \\
Photochemical & Light-harvesting materials can raise usable throughput & Show photochemical input cannot raise usable redox throughput in plausible settings \\
Material-specific & Chlorophyll/melanin are examples, not necessities & Replace them with other materials preserving the same function \\
\bottomrule
\end{tabular}
\end{center}

\section{Conclusion}
Photochemical expansion broadens the energy-input landscape; it is not the first step required for life. The release argument is a dependency graph of functions. Capture precedes overload stabilization, and modern materials are illustrative implementations rather than required historical molecules.

\section{Reproducibility Note}
The companion script \texttt{supplementary\_ool\_11.py} writes \texttt{outputs/paper\_11/photochemical\_capture\_before\_buffering.png} plus CSV/JSON summaries. No notebook is required; the static figure and small CSV/JSON files are sufficient.
\bibliographystyle{plain}
\bibliography{ool_references}
\end{document}
